\chapter{Competitive Programming Problems}
\label{chap:cpproblems}
{\huge$\bigstar$} a {\huge$\bigstar$}{\huge$\bigstar$}{\huge$\bigstar$}{\huge$\bigstar$}

\section{Organización}

El repositorio de LisPy incluye 120 problemas de programación competitiva
en \texttt{docs/problems/problem\_XXX.lp}. Cada archivo es un programa
LisPy autoejecutable. En conjunto forman una suite de regresión que verifica
el funcionamiento correcto del intérprete.

Para ejecutar todos los problemas:

\begin{lstlisting}[style=python,title=Terminal]
for %f in (docs/problems/problem_*.lp) do bin\lispy.exe %f
\end{lstlisting}

\section{Catálogo Completo}

\begin{center}
\footnotesize
\begin{tabular}{clll}
\toprule
Problema & Categoría & Dificultad & Archivo \\
\midrule
001--009 & Aritmética y números & Principiante & problem\_001.lp \\
010--018 & Listas y cons cells & Principiante & problem\_010.lp \\
019--025 & Secuencias y vectores & Intermedio & problem\_019.lp \\
026--032 & Strings y caracteres & Intermedio & problem\_026.lp \\
033--038 & Hash tables y conjuntos & Intermedio & problem\_033.lp \\
039--047 & Macros y metaprogramación & Avanzado & problem\_039.lp \\
048--055 & CLOS y POO & Avanzado & problem\_048.lp \\
056--062 & Sistema de condiciones & Intermedio & problem\_056.lp \\
063--070 & Iteración y control & Intermedio & problem\_063.lp \\
071--076 & Archivos y streams & Intermedio & problem\_071.lp \\
077--085 & Python FFI bridge & Avanzado & problem\_077.lp \\
086--090 & Paquetes y símbolos & Avanzado & problem\_086.lp \\
091--095 & Compilación y rendimiento & Experto & problem\_091.lp \\
096--100 & Misceláneos & Experto & problem\_096.lp \\
101 & Disjoint Set Union (DSU) & Experto & \texttt{problem\_101.lp} \\
102 & Pipeline middleware & Experto & \texttt{problem\_102.lp} \\
103 & Convex Hull (Graham Scan) & Experto & \texttt{problem\_103.lp} \\
104 & Weighted Task Scheduler & Experto & \texttt{problem\_104.lp} \\
105 & BUG-013 verification & Experto & \texttt{problem\_105.lp} \\
106 & Stack machine VM & Experto & \texttt{problem\_106.lp} \\
107 & Analytics engine & Experto & \texttt{problem\_107.lp} \\
108 & Delegating calculator & Experto & \texttt{problem\_108.lp} \\
109 & Cryptographic audit & Experto & \texttt{problem\_109.lp} \\
110 & Reinitialization bridge & Experto & \texttt{problem\_110.lp} \\
111 & Expression evaluator & Experto & \texttt{problem\_111.lp} \\
112 & Trading engine + Números Bouncy & Experto & \texttt{problem\_112.lp} \\
113 & Property graph engine & Experto & \texttt{problem\_113.lp} \\
114 & Introspection engine & Experto & \texttt{problem\_114.lp} \\
115 & List ops + reader macros & Experto & \texttt{problem\_115.lp} \\
116 & Reader macro DSL & Experto & \texttt{problem\_116.lp} \\
117 & CLOS MOP deep probe & Experto & \texttt{problem\_117.lp} \\
118 & Stream composition & Experto & \texttt{problem\_118.lp} \\
119 & Dynamic method protocol & Experto & \texttt{problem\_119.lp} \\
120 & Stdlib audit \& precision & Experto & \texttt{problem\_120.lp} \\
\bottomrule
\end{tabular}
\end{center}

\section{Problemas Destacados}

\subsection*{001: Precisión Racional}

Aritmética exacta con números racionales:

\begin{lispcode}
(print (+ 1/3 1/7))
\end{lispcode}

Salida: \texttt{10/21}

\subsection*{002: Fibonacci con Big Integers}

Secuencia de Fibonacci con precisión arbitraria:

\begin{lispcode}
(defun fib (n)
  (if (<= n 1) n
      (+ (fib (- n 1)) (fib (- n 2)))))
(print (fib 100))
\end{lispcode}

Salida: \texttt{354224848179261915075}

\subsection*{007: Máximo Común Divisor}

\begin{lispcode}
(print (gcd 123456789 987654321))
\end{lispcode}

Salida: \texttt{9}

\subsection*{034: Two-Sum}

Búsqueda de pares con hash table:

\begin{lispcode}
(defun two-sum (nums target)
  (let ((seen (make-hash-table)))
    (loop for n in nums for i from 0 do
      (let ((complement (- target n)))
        (multiple-value-bind (val found) (gethash complement seen)
          (if found (return (list val i))
              (setf (gethash n seen) i)))))))
(print (two-sum '(2 7 11 15) 9))
\end{lispcode}

Salida: \texttt{(0 1)}

\subsection*{048: Jerarquía de Figuras (CLOS)}

Herencia y métodos genéricos:

\begin{lispcode}
(defclass shape () ((color :initarg :color :accessor shape-color)))
(defclass circle (shape) ((radius :initarg :radius :accessor circle-radius)))
(defclass rectangle (shape) ((width :initarg :width :accessor rect-width)
                             (height :initarg :height :accessor rect-height)))
(defgeneric area (shape))
(defmethod area ((c circle))
  (* 3.141592653589793 (expt (circle-radius c) 2)))
(defmethod area ((r rectangle))
  (* (rect-width r) (rect-height r)))
(let ((c (make-instance 'circle :color "red" :radius 5)))
  (print (area c)))
\end{lispcode}

Salida: \texttt{78.53981633974483}

\subsection*{050: Especializadores EQL}

\textcolor{lisppygreen}{$\checkmark$} Corregido desde BUG-001

\begin{lispcode}
(defgeneric tax-rate (country))
(defmethod tax-rate ((country (eql :us))) 0.15)
(defmethod tax-rate ((country (eql :uk))) 0.20)
(defmethod tax-rate (country) 0.25)
(print (tax-rate :us))
\end{lispcode}

Salida: \texttt{0.15}

\subsection*{077: Python math.sqrt desde Lisp}

\begin{lispcode}
(let ((math (py:import "math")))
  (print (py:call (py:attr math "sqrt") 2)))
\end{lispcode}

Salida: \texttt{1.4142135623730951}

\subsection*{101: Disjoint Set Union (DSU)}

Implementación completa de union-find con path compression y union by rank:

\begin{lispcode}
(defclass dsu ()
  ((parent :accessor dsu-parent :initform nil)
   (rank   :accessor dsu-rank   :initform nil)
   (size   :accessor dsu-size   :initarg :size :initform 0)
   (count  :accessor dsu-count  :initform 0)))

(defun make-dsu (n)
  (let ((d (make-instance 'dsu :size n)))
    (setf (dsu-parent d) (make-array n :initial-element 0))
    (setf (dsu-rank d)   (make-array n :initial-element 0))
    (setf (dsu-count d)  n)
    (loop for i from 0 to (1- n) do
      (setf (aref (dsu-parent d) i) i))
    d))

(defmethod find ((d dsu) x)
  (let ((parent (dsu-parent d))
        (px (aref (dsu-parent d) x)))
    (if (/= px x)
        (setf (aref parent x) (find d px))
        x)))

(defmethod union ((d dsu) a b)
  (let* ((ra (find d a))
         (rb (find d b)))
    (if (= ra rb) nil
        (let ((parent (dsu-parent d))
              (rank (dsu-rank d)))
          (let ((ra-rank (aref rank ra))
                (rb-rank (aref rank rb)))
            (if (< ra-rank rb-rank)
                (setf (aref parent ra) rb)
                (if (> ra-rank rb-rank)
                    (setf (aref parent rb) ra)
                    (progn
                      (setf (aref parent rb) ra)
                      (setf (aref rank ra) (1+ (aref rank ra))))))
            (setf (dsu-count d) (1- (dsu-count d))) t)))))
\end{lispcode}

\subsection*{103: Convex Hull (Graham Scan)}

Algoritmo de geometría computacional para encontrar la envolvente convexa
de un conjunto de puntos:

\begin{lispcode}
(defclass point ()
  ((x :initarg :x :reader point-x)
   (y :initarg :y :reader point-y)))

(defun cross-product (a b c)
  (- (* (- (point-x b) (point-x a))
        (- (point-y c) (point-y a)))
     (* (- (point-y b) (point-y a))
        (- (point-x c) (point-x a)))))

(defun graham-scan (points-vec)
  (let* ((n (length points-vec))
         (pivot (find-pivot points-vec))
         (sorted (sort (copy-seq points-vec) ...))
         (hull (make-array n)) (hsize 0))
    ;; Construcción del hull con cross-product
    hull))
\end{lispcode}

\subsection*{106: Stack Machine Virtual}

Una máquina virtual completa con opcodes aritméticos, Python FFI para
funciones trascendentes, y CLOS como base:

\begin{lispcode}
(defclass vm ()
  ((stack :accessor vm-stack
          :initform (make-array 16 :fill-pointer 0 :adjustable t))
   (memory :accessor vm-memory :initform (make-hash-table :test 'equal))
   (program :accessor vm-program :initarg :program)
   (pc :accessor vm-pc :initform 0)
   (exec-count :accessor vm-exec-count :initform 0 :allocation :class)))
\end{lispcode}

\subsection*{108: call-next-method con Argumentos}

Este problema verificó la corrección de \texttt{BUG-023}, donde
\texttt{call-next-method} ignoraba los argumentos explícitos:

\begin{lispcode}
(defclass doubler (processor) ())
(defmethod compute ((p doubler) x)
  (call-next-method p (* x 2)))  ;; ahora funciona correctamente
\end{lispcode}

\subsection*{112: Números Bouncy (Project Euler)}

Un número es \textit{bouncy} si sus dígitos no son ni estrictamente
crecientes ni estrictamente decrecientes:

\begin{lispcode}
(defun bouncy-p (n)
  (let* ((s (format nil "~D" n)) (inc nil) (dec nil))
    (loop for i from 1 below (length s)
          do (when (char> (char s i) (char s (1- i))) (setf inc t))
             (when (char< (char s i) (char s (1- i))) (setf dec t)))
    (and inc dec)))

(defun euler-112 ()
  (let ((n 99) (bouncy 0))
    (loop while (/= (* bouncy 100) (* 99 n))
          do (incf n) (when (bouncy-p n) (incf bouncy)))
    n))
\end{lispcode}

\subsection*{113: Grafo con Herencia Múltiple}

Un grafo de propiedades implementado con herencia múltiple en CLOS,
listas de propiedades de símbolos, y Dijkstra:

\begin{lispcode}
(defclass property-node ()
  ((name :initarg :name :accessor pn-name :initform 'unnamed)
   (plist :initarg :plist :accessor pn-plist :initform nil)))

(defclass linked-node (property-node) ())
(defclass weighted-node (property-node) ())
(defclass composite-node (linked-node weighted-node) ())
\end{lispcode}

\subsection*{114: Metaprogramación e Introspección}

Class-of, type-of, subtypep, format condicional, equalp, coerce,
call-next-method en :before, array-rank, find-class, defstruct:

\begin{lispcode}
(print (class-of 42))     ;; \to clase INTEGER
(print (type-of "hello")) ;; \to (SIMPLE-BASE-STRING 5)
(print (subtypep 'integer 'number)) ;; \to T, T
\end{lispcode}

\subsection*{115: Operaciones de Listas y Reader/Printer}

acons, pairlis, rassoc, tailp, nconc, nreconc, adjoin, subsetp,
nsubstitute, \texttt{*print-escape*}, \texttt{*print-case*},
\texttt{*read-base*}, \texttt{copy-readtable}, \texttt{change-class},
\texttt{allocate-instance}:

\begin{lispcode}
(let ((alist (acons 'a 1 (acons 'b 2 nil))))
  (print (assoc 'a alist)))  ;; \to (A . 1)

(let* ((*read-base* 16))
  (print (read-from-string "FF")))  ;; \to 255
\end{lispcode}

\subsection*{116: Macros de Lectura Personalizadas}

Extensión del lector con \texttt{set-macro-character},
\texttt{set-dispatch-macro-character}, \texttt{make-dispatch-macro-character},
\texttt{read-delimited-list}, \texttt{*read-eval*}, \texttt{*read-suppress*}:

\begin{lispcode}
(let* ((rt (copy-readtable)))
  (set-macro-character #\$
    (lambda (s c) (declare (ignore s c)) :dollar) nil rt)
  (let ((*readtable* rt))
    (print (read-from-string "$"))))  ;; \to :DOLLAR
\end{lispcode}

\subsection*{117: CLOS Metaobject Protocol}

Introspección profunda con class-slots, slot-definition-name,
generic-function-methods, method-specializers, find-method,
ensure-generic-function, compute-applicable-methods:

\begin{lispcode}
(defclass mop-base ((name :initarg :name :accessor mb-name)))
(let ((slots (class-slots (find-class 'mop-base))))
  (print (mapcar #'slot-definition-name slots)))
\end{lispcode}

\subsection*{118: Composición de Streams}

String streams, broadcast, concatenated, echo, two-way, synonym,
file I/O, restarts, condiciones:

\begin{lispcode}
(let* ((s1 (make-string-output-stream))
       (s2 (make-string-output-stream))
       (b (make-broadcast-stream s1 s2)))
  (write-string "test" b)
  (print (list (get-output-stream-string s1)
               (get-output-stream-string s2))))
\end{lispcode}

\subsection*{120: Biblioteca Estándar y Precisión}

Auditoría completa de la biblioteca estándar con pruebas de
\texttt{make-sequence}, \texttt{fill}, \texttt{replace},
\texttt{map-into}, \texttt{mismatch}, \texttt{stable-sort},
\texttt{merge}, \texttt{vector-push-extend}, \texttt{adjust-array},
\texttt{get-universal-time}, \texttt{decode-universal-time},
\texttt{sleep}, \texttt{defsetf}, \texttt{define-setf-expander},
\texttt{apropos}, \texttt{probe-file}, \texttt{directory},
\texttt{delete-file}, y Python decimal de alta precisión.

\section{Uso como Suite de Regresión}

Estos problemas sirven como suite de regresión para el intérprete LisPy.
Cada archivo \texttt{.lp} es autoejecutable y reporta resultados al
finalizar. Los problemas 101--120 en particular documentan bugs encontrados
y corregidos durante el desarrollo:

\begin{center}
\begin{tabular}{cll}
\toprule
Bug & Descripción & Estado \\
\midrule
BUG-001 & EQL specializer no dispatcha & Corregido \\
BUG-013 & :around methods sin call-next-method & Corregido \\
BUG-015 & :before/:after requieren \allowbreak \texttt{\&allow-other-keys} & Corregido \\
BUG-016 & initialize-instance :around no dispatcha & Corregido \\
BUG-017 & with-slots falla en defmethod & Corregido \\
BUG-018 & handler-case falla en defmethod & Parcial \\
BUG-020 & multiple-value-bind en defmethod & Corregido \\
BUG-021 & Racionales a Python no convertibles & Corregido \\
BUG-022 & defmethod descarta valores múltiples & Corregido \\
BUG-023 & call-next-method ignora args explícitos & Corregido \\
BUG-025 & reinitialize-instance resetea slots & Corregido \\
BUG-026 & py:call chaining con bytes intermedios & Corregido \\
BUG-027 & shared-initialize no llamado & Corregido \\
BUG-034 a BUG-045 & 12 bugs de defmethod body & Corregido \\
\bottomrule
\end{tabular}
\end{center}

Para más detalles, consulta el archivo \texttt{BUGS.md} en la raíz del
proyecto.

\endinput
